\section{Cross-flight normalization and source-scope audit}
\label{sec:xf-audit}
The physical order used in Section~\ref{sec:xf-jacobi} is
\[
 \text{first mark foot}\ \xrightarrow{\ a\ }q_1
 \xrightarrow{\ \ell_1\ }\cdots
 \xrightarrow{\ \ell_{h-1}\ }q_h
 \xrightarrow{\ b\ }\text{second mark foot}.
\]
Only $q_1,\ldots,q_h$ are reflections. The incident collision before
the first foot and the incident collision after the second foot are
also present in the original full trajectory. Neither is deleted from
the record; neither auxiliary marked circle changes the clock.

The finite-horizon quantity $\tau_+$ denotes the fixed uniform bound
on the length of every real free flight in the original radius family,
as specified in the retained model. This revision does not certify a
new numerical value for it. All new constants display that dependency.
In particular, with
\[
 L_{\rm pair}=\tau_++97/100,\qquad
 L_{\rm end}=\tau_++191/200,
\]
$N_{\rm pair}$ and $N_\circ$ are the numbers of nonzero lattice
vectors of length at most $L_{\rm pair}$ and $L_{\rm end}$,
respectively. For either $L$, they can be evaluated exactly as
\[
 \#\{(a,b)\in\mathbb Z^2\setminus\{(0,0)\}:
                                      a^2+ab+b^2\le L^2\}.
\]
A conservative bound, without any enumeration assumption, is
$\lfloor(2L+1)^2\rfloor-1$: the disjoint open radius-$1/2$ disks
around these centers and the origin lie in the radius-$(L+1/2)$ disk;
area comparison proves the claim. This is an end-candidate count,
never a bound on the number of full physical itineraries.

The two-mark same-flight identity counts \emph{ordered} distinct
centers. For chronological pairs the expansion
$\mathscr N(\mathscr N-1)=\sum_jN_j(N_j-1)+2\sum_{j<k}N_jN_k$
accounts explicitly for both orders. A repeated center at two times
still supplies two chronological marks. The exact physical prefactor
$1/(4\pi R)$ precedes every clearance substitution, while the original
section restriction contributes $1/c_*$ only when passing to
$\nu_R^*$. Neither the factor two nor the section normalization is
absorbed into a supposed independent probability.

In the inherited identity
$E_\delta^-=\mathcal I E_\delta^+=T_RE_\delta^+$,
both expressions on the right denote \emph{set images} on the
regular invertible collision space. Thus
$\one_{E_\delta^-}(x)=\one_{E_\delta^+}(T_R^{-1}x)$ almost
everywhere; the displayed image is not $T_R^{-1}E_\delta^+$.
The null-set convention is a bulk-measure convention, not a rule that
positive one-sided endpoint traces vanish. Similarly, the critical
coarea density in Section~\ref{sec:cr-critical} is the contribution
of its stated coordinate ball and selector. It is not automatically
the complete global density near that roof value.

The new proof chain has three distinct levels. The foot standoff,
Jacobi identity, tangential-end uniqueness, exact coarea and stable
metric comparison are proved directly in
Sections~\ref{sec:xf-jacobi}--\ref{sec:xf-pressure}. The uniform
all-separation density and complete factorial bounds additionally
require the explicitly retained unconditioned physical pressure
inequality $\sup_{R,l}P_l(R)\le Q_*$. The raw pointwise theorem
still requires \eqref{eq:v112-remaining}. No conclusion at the second
level is substituted for the third. The native historical proof bodies
260--276 are preserved, and modules 277--278 alone are the new numbered
proof modules in this author revision.
